\section{Introduction}
\label{sec:introduction}

Goal-frontier maximization (GFM) proposes that an aligned agent should maximize
$\vol(G)$, the volume of the jointly achievable capability space across all
agents in a population \citep{lasser2026gfm}. The foundational paper proves
that this objective has a self-balancing property: destruction, coercion, and
self-imposed rigidity are all anti-maximizing under a single
measure-monotonicity argument. It provides a local finite-difference estimator
that reduces the requirement from volume computation to sign estimation, and
identifies conditions under which the estimator is sign-correct.

These results are conditional. They hold for any measure satisfying axioms
M1--M6 (non-negativity, null empty set, monotonicity, finite additivity,
non-triviality, superadditivity under independence), but the paper provides
only one concrete measure: the population empowerment measure $\vol_{\mathrm{PE}}$,
defined as the channel capacity of the population's joint action-state mapping.
This measure is itself intractable for the same combinatorial reasons that make
exact $\vol(G)$ computation \#P-hard \citep{dyer1988complexity}. The gap
between the theoretical framework and any implementable system is the
structure-of-$\G$ problem: what is the capability space, what measure does it
carry, and can that measure be computed?

This paper provides a constructive answer. We build $\G$ as a weighted lattice
of capabilities ordered by subsumption, discovered through two independent
channels: agent communication (what agents say they can do) and behavioral
observation (what agents are seen doing). Each capability carries a weight
derived from the information content of its benchmark. The lattice measure
$\volL$, defined via M\"{o}bius inversion, satisfies axioms M1--M6 and is
computable in polynomial time on finite lattices.

The communication boundary is the natural epistemological boundary for GFM.
Capabilities the actor cannot learn about through any channel are capabilities
it cannot reason about; treating them as zero-weight is not a limitation but an
honest accounting of the actor's epistemic position. The dual-channel design is
robust to adversarial agents: an agent that refuses to communicate its
capabilities is still observable, and an agent that overclaims is detectable
through the divergence between communicated and observed weights.

The lattice construction also upgrades the local volume estimator. On an
abstract capability space, agent reports are ternary ($R_k \in \{-1, 0, +1\}$)
because magnitudes are unavailable. On the lattice, each capability has a known
weight, so reports carry real-valued volume changes. This resolves the
acknowledged limitation that the original estimator provides no magnitude
guarantee, and strengthens the sign-correctness conditions for cases where
individual and cooperative changes have opposite signs.

\paragraph{Contributions.}
\begin{enumerate}
\item A concrete construction of $\G$ as a finite weighted poset with
  subsumption ordering (Section~\ref{sec:lattice}).
\item Dual-channel lattice discovery via communication and observation, with
  divergence as a trust signal (Section~\ref{sec:dual_channel}).
\item Proof that the lattice measure $\volL$ satisfies axioms M1--M6
  (Section~\ref{sec:axioms}).
\item A polynomial-time algorithm for $\volL(G)$ on finite lattices
  (Section~\ref{sec:tractability}).
\item An upgraded estimator with real-valued signals and magnitude bounds
  (Section~\ref{sec:estimator}).
\item A bootstrap procedure for lattice discovery from scratch
  (Section~\ref{sec:bootstrap}).
\item A worked example demonstrating self-balancing on a concrete lattice
  (Section~\ref{sec:example}).
\end{enumerate}
